\clearpage
\section*{Appendix R\quad Computational Receipts}\label{app:receipts}
\addcontentsline{toc}{section}{Appendix R: Computational Receipts}
\small\noindent Each computation marked \rcptmarker\ in the text is verified by a runnable receipt in the \texttt{receipts/} directory that ships with this work; status \textsf{OK} = verified (traced, run, output matches the claim). This appendix is generated from the receipt index, not maintained by hand.\par\medskip
\begingroup\renewcommand{\arraystretch}{1.25}
\begin{description}
\item[\rcptlabel{P03_the_two_splits}{P3R1}\textbf{P3R1}\quad\texttt{P03\_the\_two\_splits}]\hfill\textsf{[OK]}\\
\textit{sec:cubic} \ (THE DESIGNATION SPLIT AND THE SIGN SPLIT OF THE HORIZON TRIPLET AGREE EXACTLY WHEN \textbar{}r\_0\textbar{} > 1 -- A CONDITION ON THE OFFSET, NOT ON THE MASS -- AND AT NARIAI THEY NOT ONLY DIFFER BUT FAIL IN DIFFERENT PLACES. From this paper's own factorisation (r - r\_0)(r\textasciicircum{}2 + r r\_0 + r\_0\textasciicircum{}2 - 1) = 0, the designated root is r\_0 and the pair is (-r\_0 +- sqrt(4 - 3 r\_0\textasciicircum{}2))/2. ** THE SPLITS COINCIDE EXACTLY WHEN \textbar{}r\_0\textbar{} > 1, and the argument is two lines: they coincide iff the designated root is the sign-odd one, i.e. iff both pair members share a sign; the pair's product is r\_0\textasciicircum{}2 - 1 and its sum is -r\_0, so they share a sign iff r\_0\textasciicircum{}2 > 1, and that shared sign is automatically opposite to r\_0's. Checked against every sampled point of the dial. ** ** AND \textbar{}r\_0\textbar{} = 1 IS THE M = 0 MEMBER, since 2M = r\_0 - r\_0\textasciicircum{}3 -- so the boundary between agreement and disagreement is exactly the de Sitter member of the family. ** ** THE DISCRIMINATING QUANTITY IS THE OFFSET, NOT THE MASS. A first reading gave 'they differ on the positive-mass branch and agree on the negative-mass branch'; THE TABLE REFUTES IT -- agreement occurs at 2M = +0.385 and at 2M = -0.385 alike. The agreement region is INVARIANT under R : r\_0 -> -r\_0 rather than exchanged by it, R mapping the region to itself while flipping the sign of M within it. ** ON THE DIAL that is a plain angle: with r\_0 = (2/sqrt3) sin w the condition reads \textbar{}sin w\textbar{} > sqrt3/2, the sky angle exceeding 60 degrees -- and w = 60 is where sin 3w = 0, the M = 0 designation. Nariai sits at w = 30, well inside. ** AT NARIAI THE SPLITS DIFFER AND FAIL IN DIFFERENT PLACES, which is sharper than the question asked for: the designated root merges with a pair member there (asserted), so the designation split has NO CONTENT, while the sign split stays perfectly sharp isolating -2/sqrt3. A structure that survives Nariai and one that does not are not two readings of one thing. **). 
\emph{Computes:} both splits tabulated at eleven points across the dial; the \textbar{}r\_0\textbar{} > 1 predicate asserted against every sampled point; the Nariai merger asserted
 \emph{Bound:} ** AND A CAUTION THAT BELONGS WITH IT: the matter sector now uses a THIRD 1+2 split -- own wall against the other two, verified twice from independent directions. Whether that is a third structure or one of these two read on the hinges is NOT settled here and should not be assumed either way; this corpus has a standing habit of conflating threes **
\ \emph{Run:} \texttt{python3 receipts/P03\_SdS\_slicing/P03\_the\_two\_splits.py}
\item[\rcptlabel{P03_charge_obstructs_the_beads_closure_and_not_the_cosmogenesis}{P3R2}\textbf{P3R2}\quad\texttt{P03\_charge\_obstructs\_the\_beads\_closure\_and\_not\_the\_cosmogenesis}]\hfill\textsf{[ALL PASS]}\\
\textit{sec:lap} \ (**CHARGE OBSTRUCTS THE BEAD'S CLOSURE AND NOT THE COSMOGENESIS.** Two claims were being carried as one. The closed loop is obstructed TOTALLY at any nonzero charge -- the sign of f at the origin switches to +infinity -- **and the limit is singular**: the inner turning point r\_- = M - sqrt(M\textasciicircum{}2-Q\textasciicircum{}2) -> Q\textasciicircum{}2/2M shrinks to zero with the charge while the sign does NOT follow it, so the neutral case is not recovered perturbatively. The cosmogenesis theorem is untouched: its hypotheses (horizon causal structure, forced foliation, foliation-preserving morphism) mention no charge, and a sub-extremal charged collapse still carries the horizon F1 requires.). 
\emph{Computes:} r4243
 \emph{Bound:} 61
\ \emph{Run:} \texttt{python3 receipts/P03\_SdS\_slicing/P03\_charge\_obstructs\_the\_beads\_closure\_and\_not\_the\_cosmogenesis.py}
\item[\rcptlabel{P03_the_sixth_equivalence}{P3R3}\textbf{P3R3}\quad\texttt{P03\_the\_sixth\_equivalence}]\hfill\textsf{[OK]}\\
\textit{sec:hinge-geometry} \ (THE CAUSAL TRICHOTOMY RUN ON THE EXCENTRE SET: THE SET IS NULL-INERT BUT NOT CAUSALLY INERT, THE CHARACTER READS THE OWN-WALL RELATION, AND THE COMPUTATION YIELDS A SIXTH EQUIVALENCE FOR THE HINGE DISTANCE -- THE FIRST PHRASED IN THE SUBSTRATE'S OWN CAUSAL LANGUAGE RATHER THAN IN CLASSICAL TRIANGLE GEOMETRY. THE SET, DERIVED RATHER THAN QUOTED: for an equilateral triangle the excentres lie at twice the circumradius, so the excentre lines stand at transverse radius 4 alpha opposite each vertex and pierce the substrate at X\_0 = +- sqrt(15) alpha -- and they ARE substrate points (asserted), six punctures, 3 lines x 2 horns, exactly as the hinges are. ** THE EXCENTRES STAND OVER THE WALL AZIMUTHS: excentre-opposite-the-vertex and wall-antipodal-to-the-hinge give the same three azimuths, so excentre, wall and hinge are three points of one ray at radii 4 alpha, alpha and (opposite) 2 alpha. Neither fact is new; putting them side by side is. ** ** THE SET IS NULL-INERT BUT NOT CAUSALLY INERT, WHICH IS THE OPPOSITE OF WHAT WAS EXPECTED: across the 36 hinge-excentre pairs an excentre over the hinge's OWN wall is spacelike in BOTH horns while an excentre over another's is horn-dependent -- equivalently, a hinge-end and an excentre are timelike separated iff they are on opposite horns AND the excentre is not over that hinge's own wall. The causal character reads the OWN-WALL relation, the same relation the null legs read from the other direction. ** AND THE 15 EXCENTRE-EXCENTRE PAIRS REPRODUCE THE HINGE PATTERN WITH ONE CLASS MOVED: same line timelike, same horn spacelike, cross pairs TIMELIKE where the hinges' are NULL -- so the null-inertness is not an absence of structure but a radius-dependent class having moved past nullity. ** WHICH LOCATES THE HINGE RADIUS: for a cross-horn pair at transverse radius rho with azimuths 120 deg apart, s\textasciicircum{}2 = 4 alpha\textasciicircum{}2 - rho\textasciicircum{}2, vanishing at exactly rho = 2 alpha. So 2 alpha is the UNIQUE transverse radius at which the cross-horn pairs of a 120-degree-separated triple are null. **). 
\emph{Computes:} the excentre punctures asserted on the substrate; sqrt(15) asserted; the excentre and wall azimuth sets asserted identical; both pair-sets tabulated by causal character with totals asserted; the cross-pair separation solved for its null radius and asserted to be 2 alpha
 \emph{Bound:} ** WHY IT IS WORTH ADDING TO A LIST OF FIVE: this paper makes a point of the NUMBER of equivalences rather than any one -- 'their number is the point rather than any one of them' -- and all five are classical triangle geometry. A sixth that speaks the substrate's own causal language is worth more to that argument than a sixth Euclidean one. AT ITS HONEST WEIGHT: the six are NOT logically independent -- one configuration produces all of them. What is new is the KIND of statement, not the fact. AND RECORDED BECAUSE IT BEARS ON HOW THIS WAS FOUND: the conclusion was DRAFTED as 'the excentres carry no relation the substrate can see' BEFORE the table was run, on the expectation that null-inertness would extend. The table refuted it **
\ \emph{Run:} \texttt{python3 receipts/P03\_SdS\_slicing/P03\_the\_sixth\_equivalence.py}
\item[\rcptlabel{one_thirty}{P3R4}\textbf{P3R4}\quad\texttt{one\_thirty}]\hfill\textsf{[OK]}\\
\textit{\S{}488/eq:tripleathinge} \ (the 30\ensuremath{^\circ}s are ONE via triple-angle at sin w=\ensuremath{\tfrac12} \ensuremath{\Rightarrow} sin 3w=1; 12=3\ensuremath{\times}4 one route). 
\emph{Computes:} triple-angle identity 3(\ensuremath{\tfrac12})\ensuremath{-}4(\ensuremath{\tfrac12})\textsuperscript{3}=1 forces Nariai from the hinge's sine; checked on P3's dial
 \emph{Bound:} minimality (n=4 vs 8) rests on Rule 2, stated (not a theorem)
\ \emph{Run:} \texttt{python3 receipts/P03\_SdS\_slicing/one\_thirty.py}
\item[\rcptlabel{euclid7_nine_point}{P3R5}\textbf{P3R5}\quad\texttt{euclid7\_nine\_point}]\hfill\textsf{[OK]}\\
\textit{\S{}488 (nine-point)} \ (the nine-point circle of the hinge triangle IS the throat (Euler R/2=\ensuremath{\alpha}); 6th independent route to 2\ensuremath{\alpha}). 
\emph{Computes:} 9 points computed explicitly on \textbar{}X\textbar{}=\ensuremath{\alpha}; falsifier sweep rejects all placements but 2\ensuremath{\alpha}; Feuerbach\ensuremath{\to}identity
 \emph{Bound:} excircle 'third scale' question named, not pursued
\ \emph{Run:} \texttt{python3 receipts/P03\_SdS\_slicing/euclid7\_nine\_point.py}
\item[\rcptlabel{alpha_alone}{P3R6}\textbf{P3R6}\quad\texttt{alpha\_alone}]\hfill\textsf{[OK]}\\
\textit{\ensuremath{\alpha}-as-gauge} \ (\ensuremath{\alpha} alone sets the construction; 2\ensuremath{\alpha} = Thales far end (output not input); anchors forced). 
\emph{Computes:} Thales circle from the hole ends at 2\ensuremath{\alpha}; anchors \ensuremath{\surd}3\ensuremath{\alpha}/3\ensuremath{\alpha}\textsuperscript{2}/60\ensuremath{^\circ} asserted; control: tangency only at R=2\ensuremath{\alpha}
 \emph{Bound:} works Daryl's claim at the weight offered
\ \emph{Run:} \texttt{python3 receipts/P03\_SdS\_slicing/alpha\_alone.py}
\item[\rcptlabel{P03_cubic_factor_ellipse_locus}{P3R7}\textbf{P3R7}\quad\texttt{P03\_cubic\_factor\_ellipse\_locus}]\hfill\textsf{[OK]}\\
\textit{prop:factor/prop:locus} \ (cubic factors on ellipse r\textsuperscript{2}+rr\ensuremath{_0}+r\ensuremath{_0}\textsuperscript{2}=1 (A\ensuremath{_2} norm form); 3 roots sum 0; ellipse eigenstructure; involution \ensuremath{\sigma}; backward-radial reflection). 
\emph{Computes:} factorisation identity, eigenvalues \ensuremath{\tfrac12}\&3/2 (semi-axes \ensuremath{\surd}2,\ensuremath{\surd}(2/3),45\ensuremath{^\circ}), \ensuremath{\sigma}\ensuremath{\circ}\ensuremath{\sigma}=id on positive-root half, \ensuremath{\sigma}(0)=1/\ensuremath{\sigma}(1)=0/fix 1/\ensuremath{\surd}3, reflection = 2M\ensuremath{\mapsto}\ensuremath{-}2M
 \emph{Bound:} control: factorisation needs 2M=r\ensuremath{_0}\ensuremath{-}r\ensuremath{_0}\textsuperscript{3}
\ \emph{Run:} \texttt{python3 receipts/P03\_SdS\_slicing/P03\_cubic\_factor\_ellipse\_locus.py}
\item[\rcptlabel{P03_triple_angle_gnomonic}{P3R8}\textbf{P3R8}\quad\texttt{P03\_triple\_angle\_gnomonic}]\hfill\textsf{[OK]}\\
\textit{prop:gnomonic/prop:triple} \ (r\ensuremath{_0}=(2/\ensuremath{\surd}3)sin w \ensuremath{\Rightarrow} 2M=(2/3\ensuremath{\surd}3)sin 3w pure triple-angle (2/\ensuremath{\surd}3 unique); Nariai w=\ensuremath{\pi}/6 \ensuremath{\Rightarrow} r\ensuremath{_0}=1/\ensuremath{\surd}3). 
\emph{Computes:} harmonic decomp r\ensuremath{_0}\ensuremath{-}r\ensuremath{_0}\textsuperscript{3}=(\ensuremath{\rho}\ensuremath{-}\ensuremath{\tfrac34}\ensuremath{\rho}\textsuperscript{3})sin w+\ensuremath{\tfrac14}\ensuremath{\rho}\textsuperscript{3} sin3w derived; \ensuremath{\rho}=2/\ensuremath{\surd}3 unique root; Nariai=\ensuremath{\sigma} fixed point; control \ensuremath{\rho}=1 residual
 \emph{Bound:} the gnomonic-projection derivation itself is geometric, not this receipt
\ \emph{Run:} \texttt{python3 receipts/P03\_SdS\_slicing/P03\_triple\_angle\_gnomonic.py}
\item[\rcptlabel{P03_seam_continuation}{P3R9}\textbf{P3R9}\quad\texttt{P03\_seam\_continuation}]\hfill\textsf{[OK]}\\
\textit{prop:flip/sec:seam} \ (seam \ensuremath{\theta}\ensuremath{\mapsto}\ensuremath{\pi}/2+i\ensuremath{\psi}: (+,+) spherical \ensuremath{\to} (\ensuremath{-},+) de Sitter, C\textsuperscript{1} at throat, signature flip auto from d\ensuremath{\theta}=i d\ensuremath{\psi}). 
\emph{Computes:} sin(\ensuremath{\pi}/2+i\ensuremath{\psi})=cosh \ensuremath{\psi}; d\ensuremath{\theta}\textsuperscript{2}\ensuremath{\to}\ensuremath{-}d\ensuremath{\psi}\textsuperscript{2}; C0/C1 match (r=1,dr=0), C2 flips \ensuremath{-}1\ensuremath{\leftrightarrow}+1; control: real shift no flip
 \emph{Bound:} the 2D surface signature flips, not the spacetime (paper caveat)
\ \emph{Run:} \texttt{python3 receipts/P03\_SdS\_slicing/P03\_seam\_continuation.py}
\item[\rcptlabel{P03_cubic_spacing}{P3R10}\textbf{P3R10}\quad\texttt{P03\_cubic\_spacing}]\hfill\textsf{[OK]}\\
\textit{sec:temporal-threeness (spacing)} \ (the flat term decides which cubic may degenerate: disc H = -4a\textasciicircum{}4(27M\textasciicircum{}2-a\textasciicircum{}2) vanishes at M=a/3sqrt3 giving (rho,rho,-2rho) and the 2:1, while disc T = -108M\textasciicircum{}2a\textasciicircum{}4 is strictly negative and never vanishes, so T's triple is always equilateral and contributes a lone point; and in the fig-E phase the horizon roots sit at +2pi/3, 0, -4pi/3 while the turnaround sits at -2pi.cbrt(2)/3 \textasciitilde{} -151.2 deg, off the lattice by the cbrt(2) signature). 
\emph{Computes:} both discriminants factored; H's Nariai double-root factorisation verified; T's three roots distinct with equal modulus; all four phase values checked against their closed forms
 \emph{Bound:} NOT re-derived here: H-T=-a\textasciicircum{}2 r and the r=0 meeting (P7 sec:twocritical) and the four loci's derivative signatures (P7 triptych) were already placed
\ \emph{Run:} \texttt{python3 receipts/P03\_SdS\_slicing/P03\_cubic\_spacing.py}
\item[\rcptlabel{P03_turnaround_temporal_threeness}{P3R11}\textbf{P3R11}\quad\texttt{P03\_turnaround\_temporal\_threeness}]\hfill\textsf{[OK]}\\
\textit{sec:temporal-threeness} \ (the turnaround cubic's Z/3 IS the deck action of cosmic time's imaginary period 2.pi.i.alpha/3; its HALF period carries sinh\textasciicircum{}2 -> -cosh\textasciicircum{}2 (real expanding leg -> conjugate collapse law), so the wings at Im tau = +/- pi.alpha/3 are half-period images and the turnaround IS the half-period point; the temporal Z/3 factors through a branch-exchanging half-period internal to its own generator, while the spatial S\_3's exchange is R, acting between the two Nariai members rather than inside one member's root set (refined r1665: an asymmetry of where the exchange sits, not of whether one exists)). 
\emph{Computes:} r\textasciicircum{}3 invariance under the full period and the e\textasciicircum{}(2.pi.i/3) sheet cycle; sinh\textasciicircum{}2 -> -cosh\textasciicircum{}2 under the half period; sinh\textasciicircum{}2 = -1, r\textasciicircum{}3 = -2.M.alpha\textasciicircum{}2 and zero velocity at tau = -i.pi.alpha/3, with tau/period = -1/2 exactly; three real distinct horizon roots sub-Nariai
 \emph{Bound:} stated without being claimed: no identification of the two cubics' root sets (order3\_bridge.py holds that open) and no su(3) weight-triangle reading (struck)
\ \emph{Run:} \texttt{python3 receipts/P03\_SdS\_slicing/P03\_turnaround\_temporal\_threeness.py}
\item[\rcptlabel{P03_kg_is_fprime}{P3R13}\textbf{P3R13}\quad\texttt{P03\_kg\_is\_fprime}]\hfill\textsf{[OK]}\\
\textit{sec:curvature (K\_G = -f'/2r)} \ (the slicing surface's Gaussian curvature and the metric function's derivative are proportional, K\_G = -f'/(2r) identically, and the factor never vanishes for r>0 -- so the flat locus of the slice and the stationary locus of f are ONE locus described two ways, and the Hubble-Eddington radius IS the lap's rate-minimum for every member). 
\emph{Computes:} the identity asserted (K\_G + f'/2r = 0); the factor shown to have no zero for r>0; the shared zero solved and matched to (M alpha\textasciicircum{}2)\textasciicircum{}(1/3)
 \emph{Bound:} bound: an identity between two expressions, not a new physical claim; independent of the r1629 f(r\_HE)=0-iff-Nariai result
\ \emph{Run:} \texttt{python3 receipts/P03\_SdS\_slicing/P03\_kg\_is\_fprime.py}
\item[\rcptlabel{P03_rstar_merged_horizon}{P3R14}\textbf{P3R14}\quad\texttt{P03\_rstar\_merged\_horizon}]\hfill\textsf{[OK]}\\
\textit{sec:curvature (r\_star at Nariai)} \ (the Hubble-Eddington radius lies ON a horizon for exactly ONE member of the family -- the double-root one -- and there it IS the merged horizon alpha/sqrt3; the mechanism is that K\_G = 0 absorbs the mass from f, leaving the pure de Sitter form 1 - 3r\textasciicircum{}2/alpha\textasciicircum{}2, so f(r\_star)=0 iff 27M\textasciicircum{}2=alpha\textasciicircum{}2; and r\_star = A*2\textasciicircum{}(-1/3) is also the bead's unique inflection at unit speed). 
\emph{Computes:} K\_G(r\_star)=0 symbolically; f on the K\_G=0 locus reduced to the de Sitter form; f(r\_star)=0 solved for M and matched against the discriminant root independently; the Nariai double-root factorisation; f(r\_star) sign exhibited sub-, at- and over-Nariai; front-seam identity A*2\textasciicircum{}(-1/3)=r\_star
 \emph{Bound:} stated without being claimed: nothing here bears on WHY Nariai is selected -- that is P7's trichotomy
\ \emph{Run:} \texttt{python3 receipts/P03\_SdS\_slicing/P03\_rstar\_merged\_horizon.py}
\item[\rcptlabel{P03_curvature_signflip}{P3R15}\textbf{P3R15}\quad\texttt{P03\_curvature\_signflip}]\hfill\textsf{[OK]}\\
\textit{sec:curvature} \ (K\_G=1/\ensuremath{\alpha}\textsuperscript{2}\ensuremath{-}M/r\textsuperscript{3} finite+horizon-blind, one sign flip at r\ensuremath{\star}=(M\ensuremath{\alpha}\textsuperscript{2})\textasciicircum{}\ensuremath{\tfrac13}; MS/Komar return M). 
\emph{Computes:} K\_G derived via surface-of-revolution \ensuremath{-}f\ensuremath{{}^\prime}/(2r); zero at r\ensuremath{\star}; strictly increasing (one crossing); finite at horizons
 \emph{Bound:} 2D surface curvature; r\ensuremath{\to}0 divergence is areal artefact (P2)
\ \emph{Run:} \texttt{python3 receipts/P03\_SdS\_slicing/P03\_curvature\_signflip.py}
\item[\rcptlabel{P03_overcritical}{P3R16}\textbf{P3R16}\quad\texttt{P03\_overcritical}]\hfill\textsf{[OK]}\\
\textit{sec:overcritical} \ (overcritical = seam's sin\ensuremath{\to}cosh continuation on horizon angle; 2 roots complex-conjugate on ellipse's complex extension; lone real root negative; f<0 \ensuremath{\forall}r>0). 
\emph{Computes:} 3w=\ensuremath{\pi}/2+i\ensuremath{\beta}\ensuremath{\Rightarrow}sin3w=cosh\ensuremath{\beta}; discriminant<0\ensuremath{\Rightarrow}complex pair on ellipse ext; real root negative; f<0 sweep; control undercritical=3 real
 \emph{Bound:} conjugate-branch ontology cited (P15/16) not established
\ \emph{Run:} \texttt{python3 receipts/P03\_SdS\_slicing/P03\_overcritical.py}
\item[\rcptlabel{P03_charge_parity}{P3R17}\textbf{P3R17}\quad\texttt{P03\_charge\_parity}]\hfill\textsf{[OK]}\\
\textit{sec:charge} \ (RNdS slicing under mass-reflection R: R-even = 1-r\textasciicircum{}2/alpha\textasciicircum{}2+Q\textasciicircum{}2/r\textasciicircum{}2 (dS+charge), R-odd = -2M/r (mass alone); Q->-Q leaves f identical (Q via Q\textasciicircum{}2); charged horizon quartic). 
\emph{Computes:} R-parity decomposition; Q\textasciicircum{}2-only dependence; Cauchy-horizon limit
 \emph{Bound:} mass R-odd, charge R-even
\ \emph{Run:} \texttt{python3 receipts/P03\_SdS\_slicing/P03\_charge\_parity.py}
\item[\rcptlabel{P03_conjugacy}{P3R18}\textbf{P3R18}\quad\texttt{P03\_conjugacy}]\hfill\textsf{[OK]}\\
\textit{prop:conjugacy} \ (closed-form conjugacy chi o g1 = sigma o chi (both = (2/sqrt3)sin(pi/3-2a/3)) making throat-circle reflection g1 and cubic root-exchange sigma one involution). 
\emph{Computes:} symbolic + numeric (<1e-12); collapse 4-3chi\textasciicircum{}2=4cos\textasciicircum{}2(2a/3); wrong-scale control breaks
\ \emph{Run:} \texttt{python3 receipts/P03\_SdS\_slicing/P03\_conjugacy.py}
\item[\rcptlabel{P03_ellipse_foci}{P3R19}\textbf{P3R19}\quad\texttt{P03\_ellipse\_foci}]\hfill\textsf{[OK]}\\
\textit{sec:ellipse} \ (foci of the horizon-locus ellipse: eigenvalues 1/2,3/2; semi-axes sqrt2, sqrt(2/3); focal distance c=sqrt(a\textasciicircum{}2-b\textasciicircum{}2)=2/sqrt3 = the Nariai offset; axis ratio a/b=sqrt3 (equilateral)). 
\emph{Computes:} derives foci from the quadratic form; c=2/sqrt3=Nariai offset
\ \emph{Run:} \texttt{python3 receipts/P03\_SdS\_slicing/P03\_ellipse\_foci.py}
\item[\rcptlabel{P03_the_foci_are_chart_objects}{P3R20}\textbf{P3R20}\quad\texttt{P03\_the\_foci\_are\_chart\_objects}]\hfill\textsf{[OK]}\\
\textit{sec:ellipse} \ (THE QUALIFICATION THAT BELONGS WITH THE WORD 'METRIC' IN 'THE METRIC FOCI OF THE HORIZON LOCUS' -- AND IT PROTECTS THE READING RATHER THAN WEAKENING IT. ** A FOCUS IS NOT A PROJECTIVE INVARIANT OF A CONIC: IT IS DEFINED AGAINST A METRIC, and the one used is the (r, r\_0) chart's Euclidean metric, the two axes at equal scale and at right angles. Read with the A\_2 FORM ITSELF as the metric, Q = 1 is that form's unit sphere -- a CIRCLE of eccentricity zero whose foci COINCIDE at the centre. ** ** AND THE TWO READINGS GENUINELY DIFFER, WHICH THE WEYL ACTION SHOWS (both asserted): the three-cycle (r\_1,r\_2) -> (r\_2, -r\_1-r\_2) has M\textasciicircum{}3 = I, PRESERVES the A\_2 form, and is NOT orthogonal in the chart metric -- so the locus is Weyl-invariant as a SET while the Weyl action is not by chart-rotations, and the eccentricity the chart assigns is a property of the chart rather than of the configuration. ** ** THE NUMBER IS NOT THEREBY EMPTY: the focal distance is sqrt(a\textasciicircum{}2 - b\textasciicircum{}2) with the eigenvalue ratio 3 forced by the form, so 2/sqrt3 is a consequence of the form and stands however the locus is drawn. What is chart-dependent is that the consequence presents itself as a PAIR OF POINTS. ** ** HENCE A CROSS-PLANE IDENTIFICATION WITH THE OPERATOR PAPER'S NULL GENERATORS A, B CANNOT BE CANONICAL: A and B are null directions of the ambient, invariant under the isometry group; the foci are a point-pair defined against a chosen chart metric, and in the metric the configuration's own symmetry respects they do not exist as a pair at all. What survives is the statement the papers actually make -- ONE FORCED SCALE carrying the metric foci, the two null generators and the lap-close root together. **). 
\emph{Computes:} the A\_2 form's eigenvalues, semi-axes, axis ratio and focal distance with the equality to the Nariai offset asserted; and the Weyl three-cycle's order, its preservation of the form and its non-orthogonality in the chart, both asserted
 \emph{Bound:} ** A STATEMENT ABOUT WHICH METRIC A FOCUS IS DEFINED AGAINST. It does not dispute the number, the axis ratio or the scale convergence, all of which stand; it bounds what KIND of object the focal pair is **
\ \emph{Run:} \texttt{python3 receipts/P03\_SdS\_slicing/P03\_the\_foci\_are\_chart\_objects.py}
\item[\rcptlabel{P03_the_charged_collapse_does_not_form_the_eternal_inner_horizon}{P3R21}\textbf{P3R21}\quad\texttt{P03\_the\_charged\_collapse\_does\_not\_form\_the\_eternal\_inner\_horizon}]\hfill\textsf{[OK]}\\
\textit{`PO-25`; `sec:charge`, `sec:lap`} \ (the ETERNAL Cauchy horizon is NOT formed by a charged collapse of the corpus's own progenitor, in BOTH of the charge readings the row carries: the strong-cosmic-censorship ratio \$\textbackslash\{\}beta=\textbackslash\{\}alpha/\textbackslash\{\}kappa\_-\$ sits below its \$1/2\$ threshold by 250 orders (intensive) and 14 (extensive), so the obstruction `PO-25` records is a property of a STATIONARY solution the dynamical problem does not reach). 
\emph{Computes:} \$r\_-\$ and \$\textbackslash\{\}kappa\_-\$ on \$M=2.33\textbackslash\{\}times10\textasciicircum{}\{23\}M\_\textbackslash\{\}odot\$ with a deliberately GENEROUS ceiling \$\textbackslash\{\}alpha\textbackslash\{\}le1/4M\$, so the verdict is an upper bound on the horizon's chances; reproduces the row's own \$10\textasciicircum{}\{19\}\$ m and \$10\textasciicircum{}\{-64\}\textbackslash\{\},\textbackslash\{\}ell\_P\$; CONTROL at \$Q/M=0.9999\$ returns \$\textbackslash\{\}beta=17.2>1/2\$, so the criterion is non-vacuous
 \emph{Bound:} bounded: SCC holding says the eternal structure is not REACHED. It does NOT deliver a spacelike \$r=0\$, so "the branch point survives dynamically" does not follow -- what replaces the horizon is a numerical-relativity question and is not computed
\ \emph{Run:} \texttt{python3 receipts/P03\_SdS\_slicing/P03\_the\_charged\_collapse\_does\_not\_form\_the\_eternal\_inner\_horizon.py}
\item[\rcptlabel{P03_which_mass_sets_the_hubble_eddington_radius}{P3R22}\textbf{P3R22}\quad\texttt{P03\_which\_mass\_sets\_the\_hubble\_eddington\_radius}]\hfill\textsf{[OK]}\\
\textit{`PO-36`; `sec:curvature`, and `P15`'s "one constant read at two ranges"} \ (the mass question is worth a factor \$f\_b\textasciicircum{}\{-1/3\}=1.85\$ in the radius and \$1/f\_b=6.37\$ in the \$\textbackslash\{\}Lambda\$ inferred from an observed one --- and **it measures the dark fraction rather than separating the frameworks**, since \$m(r)\$ is the bend in both and whatever gravitates is in it). 
\emph{Computes:} \$r\_\{\textbackslash\{\}rm HE\}\$ on cluster, Coma-like and group scales at \$H\_0=73\$; the ratio is exactly \$f\_b\textasciicircum{}\{-1/3\}\$ and mass-independent, asserted; CONTROL at \$f\_b=1\$ gives ratio \$1.000\$ and the test falls silent, so the effect measured IS the dark fraction
 \emph{Bound:} bounded: the discrimination is quantified and its target identified; **whether a real structure's radius tracks either mass is an observation and is not made here**
\ \emph{Run:} \texttt{python3 receipts/P03\_SdS\_slicing/P03\_which\_mass\_sets\_the\_hubble\_eddington\_radius.py}
\item[\rcptlabel{P03_the_sign_switch_is_conditional_on_the_interior_mass_function}{P3R23}\textbf{P3R23}\quad\texttt{P03\_the\_sign\_switch\_is\_conditional\_on\_the\_interior\_mass\_function}]\hfill\textsf{[OK]}\\
\textit{`PO-25`'s remaining question; `sec:charge`'s obstructed/untouched separation} \ (**`P03` makes the stationary-versus-dynamical distinction once and does not apply it to the claim standing beside it.** `r6405` found the inner-horizon obstruction "a property of a stationary solution that the dynamical problem does not reach" --- and the **sign-switch** obstruction ("at any \$Q>0\$ the sign at the origin has switched") is read off the same stationary \$f\$, where \$M\$ is **constant**. \ensuremath{\Rightarrow} Dynamically **half survives**: Gauss's law makes the \$Q\textasciicircum{}\{2\}/r\textasciicircum{}\{2\}\$ term real, ***charge does not wait for stationarity to gravitate***. And **half is a condition**: the switch holds iff \$2m/r\$ stays below \$Q\textasciicircum{}\{2\}/r\textasciicircum{}\{2\}\$, ***iff \$m(r)=o(1/r)\$ as \$r\textbackslash\{\}to0\$***). 
\emph{Computes:} the stationary limit with \$M\$ constant, the competition at \$m\textbackslash\{\}sim kr\textasciicircum{}\{-p\}\$ for \$p=0,\textbackslash\{\}tfrac12,2\$, and the marginal \$p=1\$ decided by whether \$2k\$ exceeds \$Q\textasciicircum{}\{2\}\$
 \emph{Bound:}  **not claimed that \$m(r)\$ fails the condition** --- nothing here computes the interior, and mass inflation is a statement about the Cauchy horizon at finite \$r\$, not about \$r\textbackslash\{\}to0\$
\ \emph{Run:} \texttt{python3 receipts/P03\_SdS\_slicing/P03\_the\_sign\_switch\_is\_conditional\_on\_the\_interior\_mass\_function.py}
\item[\rcptlabel{P03_the_last_two_openings_await_the_same_interior}{P3R24}\textbf{P3R24}\quad\texttt{P03\_the\_last\_two\_openings\_await\_the\_same\_interior}]\hfill\textsf{[OK]}\\
\textit{`PO-25`'s condition and `PO-31`'s remainder; the owed-work instruments} \ (**nothing else is owed** --- `check\_burndown` (15 struck, 0 open, ID space intact), `check\_unworked\_blockers`, `check\_gap\_is\_held`, `check\_deferrals\_resolve` and `check\_close\_names\_work` all clean, and `FOR\_56` item 16 is **reported not failed**, a person-gate on `56`. \ensuremath{\Rightarrow} **But `PO-25`'s condition is a MODELLING task, not a reading**: \$m(r)\$ near the origin is set by the matter, and `P08` supplies the dynamics **given content** while ***not supplying the constitutive relation***. And `PO-31`'s remainder is, in its own words, "**a modelling task awaiting a progenitor interior**" --- ***two rows, two rooms, one missing interior***). 
\emph{Computes:} the instruments' reports, what `P08` supplies against what it does not, and the law/model distinction tested in both directions
 \emph{Bound:}  **not claimed a modelling run would be worthless** --- it would settle the condition **for that matter**, which is a real answer, and `P08`'s own point is that general relativity does not supply the equation of state either. What is claimed is that it is **a run with an input, not a reading**
\ \emph{Run:} \texttt{python3 receipts/P03\_SdS\_slicing/P03\_the\_last\_two\_openings\_await\_the\_same\_interior.py}
\item[\rcptlabel{C4_gnomonic_not_conformal}{P3R25}\textbf{P3R25}\quad\texttt{C4\_gnomonic\_not\_conformal}]\hfill\textsf{[OK]}\\
\textit{prop:gnomonic} \ (C4 --- which projection does CR force, and what does it cost? gnomonic: every great circle -> a straight line, NOT conformal. stereographic: conformal, great circles -> circles.). 
\emph{Computes:} runs clean; registered from its own docstring and a live run
 \emph{Bound:} description is the receipt's own wording, condensed; not independently re-derived here
\ \emph{Run:} \texttt{python3 receipts/P03\_SdS\_slicing/C4\_gnomonic\_not\_conformal.py}
\item[\rcptlabel{D2_hexagon_klein}{P3R27}\textbf{P3R27}\quad\texttt{D2\_hexagon\_klein}]\hfill\textsf{[OK]}\\
\textit{eq:tripleathinge} \ (D2 --- THE KLEIN CORRESPONDENCE on the corpus's own line configuration. CONFIRMATION 4 --- the objects, stated before anything runs: * The corpus's doubly-ruled surface is the EQUATORIAL hyperboloid -X0\textasciicircum{}2 + X1\textasciicircum{}2 + X2\textasciicircum{}2 = alpha\textasciicircum{}2, a 2-surface in R\textasciicircum{}3. Its rulings are therefore LINES IN P\textasciicircum{}3 -- which is exactly Klein's setting. * Klein sends a line in P\textasciicircum{}3 to a point of the KLEIN QUADRIC, a signature-(3,3) quadric in P\textasciicircum{}5.). 
\emph{Computes:} runs clean; registered from its own docstring and a live run
 \emph{Bound:} description is the receipt's own wording, condensed; not independently re-derived here
\ \emph{Run:} \texttt{python3 receipts/P03\_SdS\_slicing/D2\_hexagon\_klein.py}
\item[\rcptlabel{O3_charting_distance}{P3R28}\textbf{P3R28}\quad\texttt{O3\_charting\_distance}]\hfill\textsf{[OK]}\\
\textit{eq:x0proj} \ (O3 --- is r\_obs = sqrt7/2 a charting device, and does the triple-angle survive another chart? P3: r\_0 = R sin w with image radius R = tan(theta) = 1/sqrt(r\_obs\textasciicircum{}2 - 1); R is FORCED to 2/sqrt3 by the requirement that 2M = r\_0 - r\_0\textasciicircum{}3 linearise to a pure sin 3w.). 
\emph{Computes:} runs clean; registered from its own docstring and a live run
 \emph{Bound:} description is the receipt's own wording, condensed; not independently re-derived here
\ \emph{Run:} \texttt{python3 receipts/P03\_SdS\_slicing/O3\_charting\_distance.py}
\item[\rcptlabel{Q5_nariai_on_the_locus}{P3R29}\textbf{P3R29}\quad\texttt{Q5\_nariai\_on\_the\_locus}]\hfill\textsf{[OK]}\\
\textit{cor:nariai-locus} \ (Q5 --- the pencil premise, tested. Horizon cubic r\textasciicircum{}3 - r + 2M, with 2M = r0 - r0\textasciicircum{}3.). 
\emph{Computes:} runs clean; registered from its own docstring and a live run
 \emph{Bound:} description is the receipt's own wording, condensed; not independently re-derived here
\ \emph{Run:} \texttt{python3 receipts/P03\_SdS\_slicing/Q5\_nariai\_on\_the\_locus.py}
\item[\rcptlabel{I10_I15_the_horizon_is_a_turning_point_and_a_fixed_point_and_Nariai_degenerates_both}{P3R30}\textbf{P3R30}\quad\texttt{I10\_I15\_the\_horizon\_is\_a\_turning\_point\_and\_a\_fixed\_point\_and\_Nariai\_degenerates\_both}]\hfill\textsf{[OK]}\\
\textit{sec:cubic} \ (I10+I15 --- THE HORIZON IS A TURNING POINT AND A FIXED POINT, AND NARIAI DEGENERATES BOTH. Integrable-systems field bake, probes I10 and I15 together, since they are one fact from the timelike and null sides. TIMELIKE: with V = -abs(f)/2 the slicing law is a zero-energy orbit and P03's own second-derivative formula is exactly -dV/dr, verified in BOTH sign regions so the sgn(f) is checked rather than assumed. NULL: the outgoing flow dr/dv = f/2 has equilibria exactly at the zeros of f, and its linearisation IS the surface gravity f'(r\_h)/2, giving 1/4M for Schwarzschild. AT NARIAI both degenerate together: double root, V maximum at the orbit's own energy, kappa = 0 and the fixed point non-hyperbolic, and the approach costs infinite affine parameter by a logarithmic divergence. NON-DEGENERATE CONTROL: at a simple root kappa is nonzero and the approach integral converges to 2 sqrt(2) M, so the degeneracy is measured rather than assumed.). 
\emph{Computes:} runs clean; thirteen verdicts including a simple-root control that converges where Nariai diverges
 \emph{Bound:} derived here; no parameter pinned
\ \emph{Run:} \texttt{python3 receipts/P03\_slicing\_curve/I10\_I15\_the\_horizon\_is\_a\_turning\_point\_and\_a\_fixed\_point\_and\_Nariai\_degenerates\_both.py}
\item[\rcptlabel{X2r_single_periodicity}{P3R31}\textbf{P3R31}\quad\texttt{X2r\_single\_periodicity}]\hfill\textsf{[OK]}\\
\textit{sec:temporal-threeness} \ (X2-REDONE --- why the corpus computes in closed form. CONFIRMATIONS: (1) rung/level: tau-tilde is the BEAD TIME -- 'cosmic time COMPLEXIFIED... the same cosmic time continued off the real axis, which is what gives it a period' (section 0). The layer's own clock, rung 4. (3) apparatus: WHICH CUBIC? P3 separates two: H = r\textasciicircum{}3 - alpha\textasciicircum{}2 r + 2M alpha\textasciicircum{}2 the HORIZON cubic (the slicing turns;). 
\emph{Computes:} runs clean; registered from its own docstring and a live run
 \emph{Bound:} description is the receipt's own wording, condensed; not independently re-derived here
\ \emph{Run:} \texttt{python3 receipts/P03\_SdS\_slicing/X2r\_single\_periodicity.py}
\item[\rcptlabel{X3r_reality_lines}{P3R32}\textbf{P3R32}\quad\texttt{X3r\_reality\_lines}]\hfill\textsf{[OK]}\\
\textit{rem:reality-lines} \ (X3-REDONE --- what the seam actually is, analytically. CONFIRMATION 4: state the object and ask what operation is really being performed. r = sin(theta). THIS FUNCTION IS ENTIRE. The Schwarz reflection principle EXTENDS a function defined on one side of a line across it. There is nothing here to extend -- so the principle, invoked at r1869, has no work to do.). 
\emph{Computes:} runs clean; registered from its own docstring and a live run
 \emph{Bound:} description is the receipt's own wording, condensed; not independently re-derived here
\ \emph{Run:} \texttt{python3 receipts/P03\_SdS\_slicing/X3r\_reality\_lines.py}
\item[\rcptlabel{P03_the_turnaround_figure}{P3R40}\textbf{P3R40}\quad\texttt{P03\_the\_turnaround\_figure}]\hfill\textsf{[OK]}\\
\textit{Figure fig:turnaround (sec:temporal-threeness)} \ (THE CORPUS'S OTHER THREE, DRAWN FOR THE FIRST TIME. Panel (a) the complex cosmic-time plane with its ONE imaginary period 2.pi.i.alpha/3 and no real one, the fundamental domain whose identified edges make it a cylinder, the half-period wings at Im tau = +- pi.alpha/3, and the turnaround; panel (b) the three cube-root sheets, equilateral, cycled by the deck action. VERIFIED BEFORE DRAWING: the full period leaves r\textasciicircum{}3 invariant (residual exactly 0); the half period carries sinh\textasciicircum{}2 to -cosh\textasciicircum{}2; at tau = -i.pi.alpha/3 BOTH r\textasciicircum{}3 = -2 M alpha\textasciicircum{}2 and its derivative vanish, so the turnaround is genuine and not merely labelled; and the three cube roots have ONE modulus with gaps of exactly 120 degrees. CAUGHT BUG RECORDED: a first pass compared sympy Abs() objects and reported 2 distinct moduli for three equal numbers -- a canonicalisation artefact, now compared numerically). 
\emph{Computes:} the four facts symbolically, with assertions; the figure written to corpus/figures/
 \emph{Bound:} A FIGURE AND A CONTRAST, NOT A DERIVATION: it draws what P3 sec:temporal-threeness and P03\_turnaround\_temporal\_threeness already establish, and its purpose is to make the two threes distinguishable -- hinge/colour/S\_3/equatorial against turnaround/generation/Z\_3/complex-time -- since lem:twoturnings proves no affine map identifies the two root sets
\ \emph{Run:} \texttt{python3 receipts/P03\_SdS\_slicing/P03\_the\_turnaround\_figure.py}
\item[\rcptlabel{P03_the_U3_figure}{P3R43}\textbf{P3R43}\quad\texttt{P03\_the\_U3\_figure}]\hfill\textsf{[OK]}\\
\textit{Figure fig:U3 (sec:tour)} \ (THE HINGE STRUCTURE, DRAWN ONCE AND VERIFIED BEFORE IT WAS DRAWN. Panel (a) the substrate with three hinges as LINES, six punctures (3x2), six chords = twelve legs tangent at their midpoints, three walls on the throat; panel (b) the true orthographic projection to X0=0, where the projection is two-to-one onto the triangle's sides, the throat is the incircle, and w\_k is ANTIPODAL to h\_k -- with the consequence picked out in bold: a hinge's own two legs graze at the OTHER TWO hinges' walls and never at its own. CAUGHT ERROR RECORDED: the first run classified the trichotomy with the inequality reversed and swapped timelike/spacelike (6/3 for 3/6) while leaving the null count RIGHT -- the equality case is insensitive to the error that breaks the other two, which is exactly why it was checked rather than eyeballed). 
\emph{Computes:} tangency at midpoints to 4e-16; the trichotomy on all 15 pairs agreeing pairwise with the hinge/horn labels; the antipodal wall map with own-wall exclusion; asserts and exits nonzero on failure
 \emph{Bound:} A FIGURE, NOT A DERIVATION: every fact drawn is already established in P03\_twelve\_null\_legs and P03\_hexagon\_null\_triple, and this re-runs them in floating point as an independent check of the figure's OWN coordinates. No Standard Model correspondence is drawn or implied; the u/d assignment cannot be read off it at all, transitivity forbidding an absolute one
\ \emph{Run:} \texttt{python3 receipts/P03\_SdS\_slicing/P03\_the\_U3\_figure.py}
\item[\rcptlabel{P03_locus_sweep}{P3R44}\textbf{P3R44}\quad\texttt{P03\_locus\_sweep}]\hfill\textsf{[OK]}\\
\textit{sec:tour} \ (THE LOCUS SWEEP: every descriptively-named object audited against the place it actually sits. HEADLINE: one locus -- the throat point at azimuth 120k+60, i.e. r=0 on the throat circle -- carries SEVEN names (graze point, tangency point of the projection, MIDPOINT of a triangle side, vertex of the projective dual, crossing of the skew hexagon's opposite edges, the matter sector's WALL, the r=0 BRANCH POINT), and P3 sec:tour's own sentence 'One triple, THREE faces' counts three of the seven -- among the four it omits are the three that carry the physics. The sentence is not wrong; it CLOSES a list that is open. ALSO: the throat carries three names, all established identities (hole, incircle, nine-point circle); hinge and swing axis are one line; hinge vs wall scores 1/5 and seam vs wall scores 1/5 (both on the throat but at r=alpha and r=0 -- DIFFERENT loci). A THIRD CATEGORY not anticipated: one locus carrying a name the corpus has already DISOWNED (r=0 as 'the curvature singularity', which P3 calls a perspectival artefact). FOUR CANDIDATE IDENTITIES NOBODY HAS DRAWN: the lap and the cosmogenetic bead; the Nariai crest and the hinge's own tangent line; the twelve stations and the twelve null legs; and the Hubble-Eddington radius with the seam (already stated CONDITIONALLY in P3's abstract, the only identity in the corpus written in that form)). 
\emph{Computes:} the inventory grouped by locus key; the criterion run on nine pairs; proved non-identities tabulated
 \emph{Bound:} THIS IS AN AUDIT, NOT A DERIVATION: it draws together statements the corpus already makes and flags where they were never drawn together. Every identity found was implicit in two papers each stating its half -- nothing required new geometry, it required reading two sentences at once. The unpinned objects (door, lap, bead, crest, station, pseudo-sphere) are listed with what is owed, not resolved
\ \emph{Run:} \texttt{python3 receipts/P03\_SdS\_slicing/P03\_locus\_sweep.py}
\item[\rcptlabel{P03_thirds_from_closure}{P3R47}\textbf{P3R47}\quad\texttt{P03\_thirds\_from\_closure}]\hfill\textsf{[OK]}\\
\textit{sec:lap / sec:tour} \ (THE FRACTIONAL CHARGES DERIVED FROM CLOSURE, WITH NO STANDARD MODEL INPUT. DEFINITION: the three graze points sit at 60/180/300 deg, one directly between each pair of hinges, cutting the throat into three equal arcs; a constituent's WINDING is the signed number of graze points its equatorial route crosses, in thirds of a lap. FACT: between two punctures at different hinges there are exactly TWO routes, crossing 1 and 2 graze points in opposite senses, and since 1+2=3 the two routes differ by EXACTLY ONE LAP. DERIVATION: let a triple's head carry winding x and each partner y; impose only (i) x-y=1, (ii) x+2y in Z (the upper-headed triple closes, P3 sec:lap), (iii) y+2x in Z (the lower-headed one closes, the horns being symmetric). Substituting y=x-1 the totals are 3x-2 and 3x-1, integers iff 3x is. SO x = k/3 -- THE THIRDS ARE FORCED, and the 3 is the number of constituents = hinges = graze points, one 3 entering three ways. Modulo a lap the classes are k mod 3 and there are exactly THREE: k=2 gives (+2/3,-1/3) = the quarks, k=1 gives (+1/3,-2/3) = the antiquarks (the same solution under overall sign, which is C), k=0 gives (0,-1) = triality zero, the leptons and every colour singlet. AND T: the upper-headed triple (u,d,d) totals 0 and the lower-headed (d,u,u) totals +1, so the horn swap costs EXACTLY ONE LAP -- the nucleon doublet, with three of each because there are three hinges (three colour labels of one doublet)). 
\emph{Computes:} exact rational arithmetic; the route crossings enumerated on the marked azimuths; the solution classes exhausted
 \emph{Bound:} CONDITIONAL, and the antecedent is NOT shown: the derivation says IF constituents carry windings and bound objects close THEN the thirds follow. Whether a matter field here is labelled by a ROUTE rather than a point is L-74 and is now the whole gate; whether closure is a REQUIREMENT rather than a property of the vacuum slicing is L-72. Nor is the L8.1-a weld resolved (L-67, and it reaches the dimension result), nor why the k=0 class holds two states (L-69), nor anything about the continuous SU(3)
\ \emph{Run:} \texttt{python3 receipts/P03\_SdS\_slicing/P03\_thirds\_from\_closure.py}
\item[\rcptlabel{P03_winding_and_closure}{P3R48}\textbf{P3R48}\quad\texttt{P03\_winding\_and\_closure}]\hfill\textsf{[OK]}\\
\textit{sec:lap / sec:tour} \ (CHARGE AS WINDING ON THE EQUATOR, THE CLOSURE RULE, AND THE THREE INVOLUTIONS AGAINST THE SM'S THREE DISCRETE STRUCTURES. Q(u)-Q(d) = 1 EXACTLY. Split fractional/integer: u and d BOTH have fractional part 2/3; ubar and dbar BOTH have 1/3; 2/3 + 1/3 = 1. So the fractional part distinguishes QUARK from ANTIQUARK (2/3 one way round the equator against 1/3 the other -- the 120/240 arcs cut by the three hinges, one arc read two ways) and the INTEGER part distinguishes u from d, being exactly ONE LAP. CLOSURE RULE: admit a composite only if total fractional part = 0 mod 1; tested against eleven contents and agrees 11 of 11 with the observed hadron spectrum (meson/baryon/antibaryon/tetraquark/pentaquark/dibaryon close and exist; free quark/diquark/qq-qbar/qqqq/qqq-qbar do not close and do not). Confinement becomes FAILURE TO CLOSE THE LAP. And the corpus already owns the Z\_3: P3 sec:temporal-threeness derives r -> e\textasciicircum{}\{2 pi i/3\} r from cosmic time's imaginary period, for its own reasons. INVOLUTIONS: sigma -> colour (a bound triple is forced one-per-hinge by the causal trichotomy), T -> weak isospin (the horn swap costing one lap = the W carrying one unit), R -> chirality AND species at once, which is required for a Weyl field). 
\emph{Computes:} exact rational arithmetic; the closure rule enumerated over all eleven contents
 \emph{Bound:} THE IDENTIFICATION IS NOT MADE: that a constituent's fractional charge IS its equatorial arc rather than sharing its arithmetic is stated as a proof obligation. The weld conflict this bound recorded is RESOLVED: P14 sec:whichthree reseats the generations off the hinges onto the turnaround's deck Z_3, so the hinge three is the within-state index and is not asked to be generations as well (r6527). Nor do the leptons have a built seat. All registered with refutation clauses in THE\_FERMION\_SECTOR\_GEOMETRY.md. A c54.20 objection (that transitivity forbids charge coming from the thirds) is RETRACTED here: it demanded an absolute label where the structure supplies a directed arc
\ \emph{Run:} \texttt{python3 receipts/P03\_SdS\_slicing/P03\_winding\_and\_closure.py}
\item[\rcptlabel{P03_twelve_null_legs}{P3R49}\textbf{P3R49}\quad\texttt{P03\_twelve\_null\_legs}]\hfill\textsf{[OK]}\\
\textit{sec:tour} \ (THE TWELVE NULL LEGS, THE THREE TANGENCY POINTS AND THE HANDEDNESS PAIR (Daryl's statement of the structure by the numbers, verified term by term). Through each of the 6 punctures run the substrate's TWO rulings, so 6x2 = 12 null LEGS, each grazing the hole once; two legs at a tangency point make one chord (12 legs = 6 chords), and 3 tangency points x 4 legs = 12 -- the books close FOUR ways. VERIFIED: the 3 hinges are lines PARALLEL to the hole's axis crossing X0=0 at radius 2 alpha; BOTH rulings through every puncture run on through another puncture (12 of 12, and this is NOT automatic -- the substrate has a two-parameter ruling family, so it is a property of the 2 alpha placement, itself an OUTPUT per prop:twoalpha); the 12 legs graze at only THREE points, 60/180/300 deg, each DIRECTLY BETWEEN two hinges with four legs apiece, and these are prop:twoalpha(iii)'s side midpoints; a triple's two legs graze at EXACTLY -60 and +60 from its own hinge's azimuth -- MIRROR IMAGES in that azimuth, the left and right tangents, no third option; and each tangent is a NARIAI cut, sin w = 1/2 => sin 3w = 1 with r0 = +-1/sqrt3, so a triple's two legs are the Nariai member at POSITIVE and NEGATIVE gnomonic angle, graded by r0 -> -r0 = the offset parity R (which is gamma\textasciicircum{}5 in P14 and the parity making mass odd and charge even in P3 sec:charge)). 
\emph{Computes:} rulings solved exactly from v.v=0 and Q.v=0; graze points, angles, incidences and the Nariai chain all exact in sympy
 \emph{Bound:} NO Standard Model correspondence is computed or asserted here -- the counts and angles are geometry. The reading (generations, udd/duu, chirality, C at r=0, 1/3 vs 2/3 of the equator) is registered term by term with a stated refutation for each in THE\_FERMION\_SECTOR\_GEOMETRY.md. And the u/d assignment CANNOT be read off this figure: the group is transitive on the six punctures, so an absolute assignment is forbidden, not absent
\ \emph{Run:} \texttt{python3 receipts/P03\_SdS\_slicing/P03\_twelve\_null\_legs.py}
\item[\rcptlabel{P03_colour_spacelike_isospin_timelike}{P3R50}\textbf{P3R50}\quad\texttt{P03\_colour\_spacelike\_isospin\_timelike}]\hfill\textsf{[OK]}\\
\textit{sec:tour} \ (THE CAUSAL TRICHOTOMY ON THE SIX HINGE-ENDS IS A STATEMENT ABOUT WHICH OF TWO LABELS DIFFERS, AND BOTH LABELS ARE NOW NAMED -- SO ON THIS SUBSTRATE COLOUR IS A SPACELIKE RELATION AND ISOSPIN A TIMELIKE ONE. The six ends are 3 hinges x 2 horns and the fifteen pairs split 3/6/6 by which label differs (recomputed and asserted), which is the same 3/6/6 the causal classification gives -- timelike, spacelike, null -- the classification being a COMPLETE invariant since the symmetry group is transitive on the six ends with one orbit per class. ** SUBSTITUTING THE NAMES: timelike iff same colour and opposite isospin; spacelike iff same isospin and different colour; null iff both differ. ** ** HENCE THE THREE COLOURS ARE MUTUALLY SPACELIKE -- which is what 'three walls with DISJOINT SUPPORT' says in causal language, and what a baryon requires: one constituent per hinge, all in one simultaneity -- AND A STATE'S RELATION TO ITS ISOSPIN PARTNER IS CAUSAL CONTACT. ** ** AND A NULL CHORD IS THE CHANNEL THAT CHANGES BOTH LABELS AT ONCE: since a graze point IS a wall and a wall crossing is a third of a lap, the null legs are the WINDING STEP ITSELF -- not causal contact and not a gauge identification. ** ** THE TRIANGLE IS THE QUOTIENT OF THE HINGE FIGURE BY THE HORN SWAP: the six ends cover the three hinge lines 2:1 and the six null chords cover the three triangle sides 2:1 (asserted), with T the deck transformation in both -- which the corpus already stated as 'the two rulings, told apart by horn' without calling it a covering. AND D\_6 = S\_3 x Z\_2 IS THE SYMMETRY OF THAT COVER, the direct product being exactly the statement that the deck commutes with the triangle's symmetries, i.e. that the cover is REGULAR. So 'doubly ruled' downstairs and 'two-sheeted' upstairs are one fact, and the construction's Z\_2 has three descriptions -- horn swap, second ruling, deck of the covering -- not previously known to be the same Z\_2. **). 
\emph{Computes:} the six ends, the fifteen pairs tallied by which label differs with 3/6/6 asserted, the substitution table of causal character against label, and the two 2:1 projections with the fibre sizes asserted
 \emph{Bound:} ** THE TRICHOTOMY ITSELF AND THE 2:1 PROJECTION OF THE CHORDS ARE PRIOR RESULTS, verified elsewhere. What is established here is that their two labels are colour and isospin, which makes the classification a statement about two Standard Model indices rather than about six points. NOTHING HERE DERIVES EITHER IDENTIFICATION **
\ \emph{Run:} \texttt{python3 receipts/P03\_SdS\_slicing/P03\_colour\_spacelike\_isospin\_timelike.py}
\item[\rcptlabel{P03_hexagon_null_triple}{P3R51}\textbf{P3R51}\quad\texttt{P03\_hexagon\_null\_triple}]\hfill\textsf{[OK]}\\
\textit{sec:tour} \ (THE HINGE-ENDS' CAUSAL TRICHOTOMY, AND WHY THE 2+1 IT CARRIES HAS NO ABSOLUTE ASSIGNMENT (Daryl's geometry; OBJECT CORRECTED by him at c54.19 -- the c54.18 version of this receipt read the six as six independent hinges and was wrong). THE OBJECT: the hinge is the LINE the door swings on, at transverse 2 alpha, direction TIMELIKE, hence not one of the substrate's null rulings and not contained in it -- it PIERCES the substrate once per horn at X0 = +-sqrt3 alpha. THREE hinges, TWO ends apiece; the six are 3x2. (And planes through a fixed LINE form a one-parameter pencil, which is why the construction has exactly one dial.) THE TRICHOTOMY: X.Y = alpha\textasciicircum{}2(-3 eps eps' + 4 cos(th\_a - th\_b)) classifies all fifteen pairs without remainder -- TIMELIKE iff same hinge (3 pairs, X.Y = 7), spacelike iff same horn (6, X.Y = -5), NULL iff neither (6, X.Y = alpha\textasciicircum{}2). So every end is null-joined to exactly two, both opposite horn and both on OTHER hinges; two ends of one hinge CANNOT be null-separated since the hinge is timelike, so the pair a null binding appears to exclude IS THE HINGE ITSELF. The six chords project TWO-TO-ONE onto the three triangle sides (U\_i->D\_j and U\_j->D\_i, the two rulings told apart by horn), each tangent to the throat at its midpoint. THE 2+1: from any end the causal reach splits the three HINGES as \{mine, timelike\} against \{the other two, null\}. AND ITS ASSIGNMENT IS RELATIVE, PROVED: the symmetry group D\_3 x Z\_2 (order 12) acts TRANSITIVELY on the six ends and with ONE ORBIT on each causal class, so causal character is a COMPLETE invariant of a pair and nothing singles any end out). 
\emph{Computes:} all fifteen pairs, the orbit counts and the tangency verified exactly in sympy
 \emph{Bound:} ASSERTS NOTHING ABOUT uud OR udd -- and the c54.18 row's 'honest gap' (that it fails to show the two like members are alike as u and u are) is RETRACTED as a misreading: per Daryl the structure is just as well-posed as udd as it is uud, they are complementary, and species is relative in the same way, so a geometry that singled one out absolutely would be the failure. Where such a distinction would have to come from is quoted not derived: P3 sec:charge's charge Z\_2 is INDEPENDENT of Aut(A\_2)=D\_6 and fixes every root. Registered L-59, untested
\ \emph{Run:} \texttt{python3 receipts/P03\_SdS\_slicing/P03\_hexagon\_null\_triple.py}
\item[\rcptlabel{P03_two_two_plus_ones}{P3R52}\textbf{P3R52}\quad\texttt{P03\_two\_two\_plus\_ones}]\hfill\textsf{[OK]}\\
\textit{sec:cubic / sec:ellipse} \ (**!! MARKED c54.159: PART 2's test of whether the SIGN split moves when the designation is cycled IS a loop that recomputes pos and neg from the root set alone and never uses the designated root.** Its reported invariance cannot fail, so PART 2's geometry-versus-vantage conclusion rests on a tautology rather than a test). 
\emph{Computes:} partitions computed along the dial and at both Nariai designations; the divisibility exact in sympy at D=4..8
 \emph{Bound:} NAMES the objects FIGURE R3 failed to name; it is NOT evidence for or against any mapping to uud, and asserts no match. The candidate sides it records are candidacy, not identification
\ \emph{Run:} \texttt{python3 receipts/P03\_SdS\_slicing/P03\_two\_two\_plus\_ones.py}
\item[\rcptlabel{P03_dimension_collapse}{P3R53}\textbf{P3R53}\quad\texttt{P03\_dimension\_collapse}]\hfill\textsf{[OK]}\\
\textit{sec:family / rem:dimension} \ (for which D does the horizon relation collapse to a single multiple-angle in the sky angle? 2M = r0\textasciicircum{}(D-3) - r0\textasciicircum{}(D-1): the harmonics below the top number ONE at D=4,5 and TWO OR MORE from D=6, against one available scale, so D>=6 has no pure multiple-angle and no forced fold (checked for a single harmonic of ANY order, not just the top). D=5 collapses at varrho=1 to a 4-fold; D=4 at 2/sqrt3 to prop:triple's sin 3w. TWO CHECKS THAT COULD HAVE FAILED: the (D-1)-gon hinge placement returns the FORCED 2*alpha at D=4 by a route unrelated to how P3 fixes it, and in both surviving D the hinge tangent lands on the Nariai radius computed independently from f's double root. Also the adversarial pass run BEFORE propagation -- five attacks, none lands). 
\emph{Computes:} calibrated on two forced 4D values it was not given (2*alpha, r\_N); D=4 reproduces prop:triple exactly; the D>=6 same-mass residual is O(1), not marginal
 \emph{Bound:} assumes the D-dimensional f is Tangherlini-de Sitter -- an extension of the codimension-one operator, not a rederivation; stated in the paper. Records a defect it caught in itself (an odd extension of 2M wrong at odd D)
\ \emph{Run:} \texttt{python3 receipts/P03\_SdS\_slicing/P03\_dimension\_collapse.py}
\item[\rcptlabel{P03_reach_probe_deepening}{P3R54}\textbf{P3R54}\quad\texttt{P03\_reach\_probe\_deepening}]\hfill\textsf{[OK]}\\
\textit{sec:family / rem:dimension} \ (the REACH programme's own probes re-run at general D, since every bake was run against a four-dimensional cut at a time when nothing distinguished 'structural' from 'true at D=4'. SIX PROBES, and they do not all go one way: OPTICS O1 STRENGTHENED (r\_photon = r\_Nariai in every D tested, so the tangency/photon-sphere identity is dimension-independent) - CATEGORY K8b GROUNDED (its bound is so(4,1), the stabiliser of a FOUR-dimensional cut, dim 10 in every dS\_D) - CATEGORY K5 NARROWED (the isotropy enhancement is +2 in every D; the pair (6,4) is D=4) - OPTICS O6 VERDICT WIDER THAN ITS COMPUTATION (2f - r\textasciicircum{}2 f'' is the constant 2 only at D=4) - QUADRIC Q4's self-discount too strong (a cross-ratio needs four points; the vantage set has D of them, so it is a cross-ratio configuration precisely at D=4, and the D=5 subset is harmonic not equianharmonic) - QUADRIC 4d ANSWERED from the untried queue (the hinge triangle is NOT self-polar; and a plane simplex needs three points, so the question exists only at D=4)). 
\emph{Computes:} each probe re-derived from its own computation, not annotated; O1's identity checked to exact radicals at D=4..8; the 4d negative is exact (alpha/2 vs alpha)
 \emph{Bound:} assumes the D-dimensional f is Tangherlini-de Sitter, as rem:dimension does. NONE of the six is a second vote for D=4 - all descend from the one fold, and the receipt says so
\ \emph{Run:} \texttt{python3 receipts/P03\_SdS\_slicing/P03\_reach\_probe\_deepening.py}
\item[\rcptlabel{P03_step3_sweep}{P3R55}\textbf{P3R55}\quad\texttt{P03\_step3\_sweep}]\hfill\textsf{[OK]}\\
\textit{sec:tour / rem:dimension} \ (**!! MARKED c54.159: FAMILY A's columns "tan len = height?" and "midpt tangency?" print True at every D as if verified, and the first computes `sqrt(Rv**2-1) - sqrt(Rv**2-1) == 0`, a tautology independent of Rv, while the second hardcodes `midpt = True`.** So the two equivalences the receipt calls dimension-independent and load-bearing -- prop:tangentnull and power-of-a-point = height\textasciicircum{}2 -- are the only two of the five with no computation behind them). 
\emph{Computes:} each family re-derived from the D-dimensional f, not annotated; the D=4 row reproduces the corpus's own values in every family
 \emph{Bound:} assumes the D-dimensional f is Tangherlini-de Sitter, as rem:dimension does. The families overlap and every dimension statement descends from the ONE fold -- not 49 confirmations
\ \emph{Run:} \texttt{python3 receipts/P03\_SdS\_slicing/P03\_step3\_sweep.py}
\item[\rcptlabel{P03_step3_refused}{P3R56}\textbf{P3R56}\quad\texttt{P03\_step3\_refused}]\hfill\textsf{[OK]}\\
\textit{sec:tour / rem:dimension} \ (THE THREE ITEMS AN EARLIER PASS RECORDED INSTEAD OF WORKING, worked. (I) THE ANTIPODAL Z\_2 vs THE OFFSET PARITY -- KILLED on an exact computation: the antipodal map has det +1 and REVERSES X\_0, the offset parity diag(1,1,-1,1,1,1) has det -1 and FIXES X\_0, so they cannot be one Z\_2 read at two levels; and the kill is informative because the antipodal map is exactly T composed with spatial inversion, putting it on T's side of the corpus's own T/R line -- the quadric ledger's refusal of the projective quotient is a TIME-orientation argument while the chirality-grading parity is the SPATIAL one. (II) FIGURE R3 -- A KILL WAS WRITTEN AND IS RETRACTED (same session, Daryl): P14 names TWO 2+1s in one sentence (the designation 2+1, sigma's own, and the sign 2+1) and R3 names neither, so a computation about the sign one cannot close the row; the designation 2+1 is a flavour-side structure and P14 records that the locus DOES carry a flavour-breaking structure. R3 REMAINS OPEN and wider. What IS established: the row's unsourced reason is WITHDRAWN (the split is not a real-slice artifact, since sign(r) is species), the sign 2+1 tracks D's parity exactly so it is the mass-parity in another register, and the corpus carries two 2+1s the catalogues had not distinguished. (III) the category bake's two ZEROS adjudicated in prose: they stand as verdicts about the field). 
\emph{Computes:} the involution comparison is exact matrix algebra; the root-sign counts computed at a sub-critical mass in every D
 \emph{Bound:} THE RETRACTION IS THE POINT: a narrowly-scoped result was chained with two others into an absolute foreclosure of an open search -- the cyanide face, and the second time that drill-site has been imploded and restored. The Nariai/three-quark lead and R3 both stand OPEN
\ \emph{Run:} \texttt{python3 receipts/P03\_SdS\_slicing/P03\_step3\_refused.py}
\item[\label{rcpt:C4_the_reality_grid_and_the_pencil}\textbf{}\quad\texttt{C4\_the\_reality\_grid\_and\_the\_pencil}]\hfill\textsf{[OK]}\\
\textit{sec:tour} \ (sinh-s reality set is a pencil not a grid; sinh\textasciicircum{}2-s two lowest lines are the bead-s branches). 
\emph{Computes:} Im sin=cos x sinh y vs Im sinh=sin y cosh x; Im sinh\textasciicircum{}2 lines at k pi/2; the bead bound is the second (11)
 \emph{Bound:} NOT-A-PAPER-CLAIM --- discharges L-203 probe C4
\ \emph{Run:} \texttt{python3 receipts/RM\_C\_complex\_analysis/C4\_the\_reality\_grid\_and\_the\_pencil.py}
\item[\rcptlabel{A6_item_58_resolves_split}{P3R57}\textbf{P3R57}\quad\texttt{A6\_item\_58\_resolves\_split}]\hfill\textsf{[OK]}\\
\textit{\S{}sec:hinge-geometry (excentre)} \ (ITEM 58 RESOLVES **SPLIT**, AND THE SPLIT EXPOSES THE UNBANKED SWEEP'S THIRD FAILURE MODE: A RESULT BANKED WITHOUT ITS WORKING NAME. `unbanked.py` surfaced `excentre` at 35 uses across five receipts and ZERO across all seventeen papers, so two results were routed as unbanked --- a sixth equivalence for the hinge distance \$2\textbackslash\{\}alpha\$ phrased in causal language, and the hinge configuration's null connectivity CLOSED at 0 of 36 pairs. **Registered c54.225 (`L-559`): the file existed and had NO INDEX row, so it was never run by the gate, never in an appendix, and never in the assertion census.**). 
\emph{Computes:} resolves the routed item against the papers' own text and the receipt corpus, separating the result that IS banked from the working NAME that is not
 \emph{Bound:} NOT-A-PAPER-CLAIM --- routed-item resolution. **Registered, not written, at c54.225: the row is the repair; the content is r2678's and is unaltered.**
\ \emph{Run:} \texttt{python3 receipts/P03\_SdS\_slicing/A6\_item\_58\_resolves\_split.py}
\item[\rcptlabel{AG1_the_two_discriminants_are_a_resultant_apart}{P3R58}\textbf{P3R58}\quad\texttt{AG1\_the\_two\_discriminants\_are\_a\_resultant\_apart}]\hfill\textsf{[OK]}\\
\textit{sec:ontology} \ (**THE SLICING DISCRIMINANT AND THE HORIZON CUBIC'S ARE THE TWO FACTORS OF THE CLASSICAL REDUCIBLE-CURVE DISCRIMINANT FORMULA --- AND THE CORPUS COMPUTES BOTH AND NAMES NEITHER.** The corpus states the factoring \$(r-r\_0)(r\textasciicircum{}2+rr\_0+r\_0\textasciicircum{}2-1)\$ in `P03` and `P12`. Its two degenerations are geometrically different events: the **conic's own two roots colliding**, whose discriminant in \$r\$ is exactly \$4-3r\_0\textasciicircum{}2\$ --- *which IS `P03`'s slicing discriminant, and is why `P03` reads its three regimes off it* --- and the **line meeting the conic**, \$\textbackslash\{\}mathrm\{Res\}(\textbackslash\{\}text\{line\},\textbackslash\{\}text\{conic\})=3r\_0\textasciicircum{}2-1\$. \ensuremath{\Rightarrow} ***And the whole cubic's discriminant is their product with the resultant SQUARED:*** \$\textbackslash\{\}mathrm\{disc\}(C)=\textbackslash\{\}mathrm\{disc\}(\textbackslash\{\}text\{conic\})\textbackslash\{\}cdot\textbackslash\{\}mathrm\{Res\}\textasciicircum{}2=(4-3r\_0\textasciicircum{}2)(3r\_0\textasciicircum{}2-1)\textasciicircum{}2\$, the textbook formula for a reducible plane curve.  **AND THEY ARE NOT THE SAME VANISHING CONDITION:** the cubic's discriminant has FOUR zeros in \$r\_0\$ and the slicing one has TWO --- the squared factor adds \$r\_0=\textbackslash\{\}pm1/\textbackslash\{\}sqrt3\$, where \$4-3r\_0\textasciicircum{}2=3\$ and is nowhere near zero. *So "the discriminant vanishes" is ambiguous in this corpus until you say which, and that ambiguity is exactly what r3547 repaired in `P3`.*  `ROUTED, NOT APPLIED` --- the paper-side repair is 59's and is done; what is owed is the NAME, `resultant` being \ensuremath{\times}0 across seventeen bodies). 
\emph{Computes:} COMPUTES: the reducible form asserted identical to \$r\textasciicircum{}3-r+2M\$ with \$2M=r\_0(1-r\_0\textasciicircum{}2)\$; the conic's discriminant asserted equal to `P03`'s \$4-3r\_0\textasciicircum{}2\$; the resultant asserted equal to the conic evaluated at \$r\_0\$ AND to \$3r\_0\textasciicircum{}2-1\$; the product identity \$\textbackslash\{\}mathrm\{disc\}(C)=\textbackslash\{\}mathrm\{disc\}(\textbackslash\{\}text\{conic\})\textbackslash\{\}cdot\textbackslash\{\}mathrm\{Res\}\textasciicircum{}2\$ asserted exactly; and the zero-sets asserted to be a PROPER subset relation; **and every one of the four zeros asserted to carry \$\textbar{}2M\textbar{}=2/(3\textbackslash\{\}sqrt3)\$ exactly and to leave the cubic with a repeated root, with the role of \$r\_0\$ at each asserted --- simple where the conic's own roots collide, repeated where the resultant vanishes**
 \emph{Bound:} **NOT-A-PAPER-CLAIM** --- no paper edited. **NOT claimed:** that 59's r3547 repair is wrong --- it is right, and "they vanish together at Nariai and differ by an exact square" is true at Nariai; what this adds is that the zero-sets differ away from it. **NOT claimed** that the corpus should import algebraic geometry --- only that it computes both factors of a named classical formula without naming it
\ \emph{Run:} \texttt{python3 receipts/P03\_slicing\_curve/AG1\_the\_two\_discriminants\_are\_a\_resultant\_apart.py}
\item[\rcptlabel{P03_acceleration_is_slice_curvature}{P3R12}\textbf{P3R12}\quad\texttt{P03\_acceleration\_is\_slice\_curvature}]\hfill\textsf{[OK]}\\
\textit{sec:curvature \ensuremath{\cdot} rem:twocritical \ensuremath{\cdot} the bend} \ (on ANY member of the energy family (E cancels: (dr/dtau)\textasciicircum{}2 = E\textasciicircum{}2 - f, and E is a constant of the motion) d\textasciicircum{}2r/dtau\textasciicircum{}2 = -f'/2 = r*K\_G identically, so THE SIGN OF THE SLICE'S GAUSSIAN CURVATURE IS THE SIGN OF THE COMOVING ACCELERATION -- and the flat locus is the deceleration-to-acceleration turnover, which on Nariai is the seam and in standard terms is rho\_m = 2 rho\_Lambda (csch\textasciicircum{}2 w = 2 <=> sinh w = 1/sqrt2, the seam's phase). A recent cosmological epoch: 7.06 Gyr at H0=73, 7.65 at 67.4, preceding matter-Lambda equality [borrowed from P3 \ensuremath{\cdot} P7 \ensuremath{\cdot} P8]). 
\emph{Computes:} the identity asserted; the K\_G zero solved and matched to (M alpha\textasciicircum{}2)\textasciicircum{}(1/3); csch\textasciicircum{}2 w = 2 <=> sinh w = 1/sqrt2 proved by algebra after simplify() failed on the closed forms; epochs at both H0
 \emph{Bound:} bound: identities on the Nariai E=1 worldline; no causal claim between seam and turnover -- they are one locus; dating inherits H0
\ \emph{Run:} \texttt{python3 receipts/P03\_SdS\_slicing/P03\_acceleration\_is\_slice\_curvature.py}
\item[\rcptlabel{P07_lap_timeline}{P7R5}\textbf{P7R5}\quad\texttt{P07\_lap\_timeline}]\hfill\textsf{[OK]}\\
\textit{rem:twocritical + P3 sec:temporal-threeness} \ (A2.11's timeline: the landmark phases exact, w\_back = 2 w\_front = arccosh 2 = ln(2+sqrt3) (the 2:1 in radius as a 1:2 in phase, by the identity cosh2w = 1+2sinh\textasciicircum{}2w across the sinh/cosh legs); the intervals in Gyr on the corpus's constants; and THE LIFT ADVANCES Re(tau) BY EXACTLY ZERO, so the collapse-to-expansion handoff takes no cosmic time [borrowed from P7]). 
\emph{Computes:} phases in closed form, the relation proved via the identity and cross-checked as algebra and numerically; A = 2\textasciicircum{}(1/3) rho; the Nariai condition asserted as self-consistency rather than against a number; the present epoch checked against the independent flat-LCDM age formula to 1e-12; the 1/H0 scaling exposed
 \emph{Bound:} held: tau is the FUNDAMENTAL OBSERVERS' proper time, not the exterior static clock (P1's infinite-time result is the exterior one). TWO LEVELS, never collapsed: GEOMETRICALLY the lap passes THROUGH r=0 -- a branch point, not a barrier, the substrate smooth across it, and the crossing is where the species reverses; PHYSICALLY no finite-time layer S\_t CONTAINS the curvature singularity (the Smoothness theorem, about the layers and not the curve), so the cosmological future is seeded at the finite-curvature seam. P7: 'claims at different levels, not competing ones.' The rho-crossing figure is a law-phase, not the seeding. The imaginary length is contour, not duration, and not recurrence. *r1632 collapsed the second level into the first and wrote 'r=0 is an unattained limit' -- inverting the corpus's central result; corrected*
\ \emph{Run:} \texttt{python3 receipts/P07\_CR\_framework/P07\_lap\_timeline.py}
\item[\rcptlabel{C5_sct_already_linear}{P3R26}\textbf{P3R26}\quad\texttt{C5\_sct\_already\_linear}]\hfill\textsf{[OK]}\\
\textit{rem:reality-lines} \ (TWO CLOSURES. (A) X3r's open question: sin is real on the real axis AND on every line Re z = pi/2 + k.pi. The crossings are at z = pi/2 + k.pi. Are the odd ones the corpus's BACK seam as the even ones are its FRONT? (B) The conformal queue's last entry: 'C1 computed the LINEAR conformal maps only; the SPECIAL CONFORMAL TRANSFORMATIONS are untested -- the one place the conformal field could still bite.' [borrowed from P3 (+P5)]). 
\emph{Computes:} runs clean; registered from its own docstring and a live run
 \emph{Bound:} description is the receipt's own wording, condensed; not independently re-derived here
\ \emph{Run:} \texttt{python3 receipts/P03\_SdS\_slicing/C5\_sct\_already\_linear.py}
\item[\rcptlabel{O1_photon_sphere_nariai}{P7R15}\textbf{P7R15}\quad\texttt{O1\_photon\_sphere\_nariai}]\hfill\textsf{[OK]}\\
\textit{sec:applications-synthesis} \ (O1 --- the photon sphere of the SdS family, and where it sits at Nariai. f(r) = 1 - 2M/r - r\textasciicircum{}2/alpha\textasciicircum{}2. Photon sphere: d/dr (f/r\textasciicircum{}2) = 0. Impact parameter b\textasciicircum{}2 = r\textasciicircum{}2/f. [borrowed from P7 (+P3)]). 
\emph{Computes:} runs clean; registered from its own docstring and a live run
 \emph{Bound:} description is the receipt's own wording, condensed; not independently re-derived here
\ \emph{Run:} \texttt{python3 receipts/P07\_CR\_framework/O1\_photon\_sphere\_nariai.py}
\item[\rcptlabel{P03_operator_at_general_D}{P3R36}\textbf{P3R36}\quad\texttt{P03\_operator\_at\_general\_D}]\hfill\textsf{[OK]}\\
\textit{rem:dimension / thm:kernel} \ (** ASSUMPTION (A1) DISCHARGED. ** CR's OWN operator at general D returns Tangherlini-de Sitter and nothing else. The point previously missed: CR's operator is NOT a metric function -- P8 thm:kernel defines the vacuum sector as the KERNEL OF THE MATTER FUNCTIONAL on a cut in the construction gauge ('derived, not matched'), and that condition runs at any D. Computed from the metric rather than quoted: for ds\textasciicircum{}2 = -f dt\textasciicircum{}2 + dr\textasciicircum{}2/f + r\textasciicircum{}2 dOmega\textasciicircum{}2\_\{D-2\}, G\textasciicircum{}t\_t = (D-2)/(2r\textasciicircum{}2)[r f' + (D-3)(f-1)] (verified at D=4,5,6), so T=0 is the first-order linear ODE r f' + (D-3)(f-1) + 2 Lambda r\textasciicircum{}2/(D-2) = 0 -- P8's own equation at D=4 -- whose ENTIRE solution space at D=4,5,6,7,8 is f = 1 - 2M/r\textasciicircum{}(D-3) - 2 Lambda r\textasciicircum{}2/((D-1)(D-2)) with ONE constant of integration, verified identical to Tangherlini-dS at every D tested. With alpha(D) = sqrt((D-1)(D-2)/2Lambda), which returns sqrt(3/Lambda) at four, this is f = 1 - 2M/r\textasciicircum{}(D-3) - r\textasciicircum{}2/alpha\textasciicircum{}2 at every D, and the dimension result's input 2M = r\_0\textasciicircum{}(D-3) - r\_0\textasciicircum{}(D-1) is that f at a root. THE DIMENSION RESULT THEREFORE STRENGTHENS: the collapse, the fold being D-1 and the mass parity all run on a function the construction GENERATES rather than borrows [borrowed from P3 / P8]). 
\emph{Computes:} the Einstein tensor computed from the metric at D=4,5,6 and matched to the closed form; the ODE solved by dsolve at D=4..8 and compared term by term with Tangherlini-dS
 \emph{Bound:} TWO INPUTS REMAIN AND ARE NAMED: (i) THE LOCK -- the construction gauge g\_tt g\_rr = -1 is assumed at general D, and P8 says the single slicing curve enforces it at four but this does not show it does so at general D (L-89; P8's lapse split is the machinery that would say otherwise); (ii) THE SPHERE -- the transverse space is taken round S\textasciicircum{}\{D-2\}, which at general D is a choice (L-90). And (A2), the single-harmonic collapse being a REQUIREMENT at other D, is untouched
\ \emph{Run:} \texttt{python3 receipts/P03\_SdS\_slicing/P03\_operator\_at\_general\_D.py}
\item[\rcptlabel{P03_slate_worked}{P3R39}\textbf{P3R39}\quad\texttt{P03\_slate\_worked}]\hfill\textsf{[OK]}\\
\textit{sec:alpha-action / sec:synchronous / range} \ (THE NINE INHERITED BAKE QUEUES, WORKED (L-24, L-25, L-26, L-27, L-28, L-30, L-31, L-32, L-33). HEADLINE: L-26 -- the lead's premise was WRONG and the truth is stronger. Sub-ring spacing is exp(-2 pi lambda/Omega); O5 shows lambda -> 0 at the forced member but Omega -> 0 with it, both proportional to f(r\_ph). And for f = 1 - 2M/r - r\textasciicircum{}2/alpha\textasciicircum{}2 one has the IDENTITY f'' = 2(f-1)/r\textasciicircum{}2, whence lambda\textasciicircum{}2 - Omega\textasciicircum{}2 = 0 IDENTICALLY at every radius and every member -- SO lambda/Omega = 1 EXACTLY and the sub-ring demagnification is exp(-2 pi) = 0.0018674, the Schwarzschild value, INDEPENDENT OF M AND OF Lambda and unchanged at the degenerate member; what degenerates there is the TIMESCALE, not the spacing. L-24: the confocal family of an equilateral quadric IS its own alpha-family, which the corpus carries as the R+ of O(5,1) x R+ -- so the dilation is CONFOCAL. L-25: J\textasciicircum{}+ is at infinite affine parameter and the seam at finite proper distance, and P8 sec:synchronous already says so. L-27: the comoving reading is ONE universal aberration factor sqrt((1-v)/(1+v)) with v = sqrt(1-f), not a redo. L-28: the orbit space X/G is P9's range and the stratification is P12's. L-30: ZERO null connections among excentres and ZERO excentre-to-hinge (0 of 36) -- the hinge configuration's null connectivity is CLOSED. L-31/L-33: the Euclid protocol was RUN (PART XII, r1110) and the SCTs were RUN (C5, r1888) -- three slate items are STALE QUEUE TEXT quoting the pre-run question [borrowed from P3 / P5 / P8 / P9]). 
\emph{Computes:} the identity proved symbolically; photon-sphere quantities computed at four members; excentre incidences exhausted over all 36 pairs; the ledger texts read
 \emph{Bound:} ** NOT-A-PAPER-CLAIM: this is a PROCESS receipt -- it records a batch of leads worked together, not a claim any paper makes, so it is not owed a citation. ** L-32 is a RENDERING task, not a question, and remains undrawn (L-106). L-24's orthogonality theorem is identified as unused, not used (L-99). L-30 is a computed negative -- inert under NULLITY, which does not exclude other relations (L-104)
\ \emph{Run:} \texttt{python3 receipts/P03\_SdS\_slicing/P03\_slate\_worked.py}
\item[\rcptlabel{P03_the_adjoint_is_entailed}{P3R42}\textbf{P3R42}\quad\texttt{P03\_the\_adjoint\_is\_entailed}]\hfill\textsf{[OK]}\\
\textit{prop:factor / sec:cubic} \ (THE ADJOINT OF SU(3) IS ENTAILED, NOT OWED -- AND THE CENTRE AND THE ADJOINT ARE ONE DATUM. P3's factorisation ellipse r\textasciicircum{}2 + r r\_0 + r\_0\textasciicircum{}2 = 1 is, computed rather than compared, the unit circle of the KILLING FORM on the Cartan subalgebra of su(3): Vieta puts the three horizon roots in the sum-to-zero plane, whose induced Euclidean norm is exactly TWICE the corpus's own A\_2 norm form. The six designation transitions e\_i - e\_j -- the corpus's own sigma, the step to a neighbouring hinge -- satisfy all four root-system axioms in that form, INTEGRALITY included (Cartan integers -2,-1,1,2), with angles 0/60/120/180 and no 90, and Weyl group of order 6. ALL SIX ROOTS HAVE CORPUS-NORM EXACTLY 1, so they LIE ON P3's OWN ELLIPSE: the corpus drew the A\_2 root hexagon and called it a factorisation identity. Six roots at rank two leaves only A\_2, which determines sl(3,C) uniquely, of dimension \textbar{}Phi\textbar{} + rank = 6 + 2 = 8 -- the hexad plus the sum-to-zero plane, counted. And \textbar{}Z(SU(3))\textbar{} = det(Cartan) = 3 comes from the SAME object, so one cannot have the centre of A\_2 without having A\_2 [borrowed from P3 / P13 / P14]). 
\emph{Computes:} the Killing/corpus form ratio computed symbolically; the four axioms tested on the six roots; the Weyl group generated by closure (a first pass took two-letter words and got the rotation subgroup, 3, which is recorded); the Cartan matrix and its determinant
 \emph{Bound:} ** THE BOUND IS THE LARGER HALF: a Lie algebra being present is not a gauge symmetry being present. No bundle, no connection, and no reason the content sits in the fundamental. The su(3)-not-in-so(5,1) wall (P13) and the global-Wick placement on S\textasciicircum{}5 stand untouched. What changes is the SHAPE of the debt -- the corpus owes the object su(3) acts on, not the adjoint -- which is P14's own open bundle question **
\ \emph{Run:} \texttt{python3 receipts/P03\_SdS\_slicing/P03\_the\_adjoint\_is\_entailed.py}
\item[\rcptlabel{P03_sheet_index_and_factors}{P3R45}\textbf{P3R45}\quad\texttt{P03\_sheet\_index\_and\_factors}]\hfill\textsf{[OK]}\\
\textit{sec:temporal-threeness / sec:lap / prop:wall} \ (WHAT TRANSFORMS AS A WALL-CROSSING COUNT, AND HOW THE LABELS FACTORISE. RETRACTS a prior inference that 'wall<->hinge is a bijection so they are the same three' -- INVALID, since colour and generation in the SM are both threes, both permuted, and independent. CRITERION run on wall vs hinge (LOCUS, TYPE, DEFINITION, EQUIVARIANCE, SEPARABILITY): four of five say DIFFERENT (a hinge is a timelike LINE at 2 alpha with two punctures and no legs; a wall is a POINT on the throat at alpha with no punctures and four legs), and the only 'same' is EQUIVARIANCE, which is exactly the property that does not force labels to coincide. INCIDENCE: a hinge's legs cross the other two hinges' walls and never its own, so realised (hinge,wall) pairs number 6 of 9 -- the product minus the antipodal diagonal, a constraint hence a prediction. FACTORISATION: a leg is (hinge, horn, which-of-the-other-two-walls) = 3x2x2 = 12, and the third factor is DERIVED not assigned -- the two legs graze at -60/+60, each a Nariai cut at w = -+30, r0 = -+1/sqrt3, and r0 -> -r0 is R = gamma\textasciicircum{}5. ANSWER: r is a cube root of an invariant (r\textasciicircum{}3 = 2M a\textasciicircum{}2 sinh\textasciicircum{}2(3 tau/2a)) whose Z\_3 deck action tau -> tau + 2 pi i a/3 sends r -> e\textasciicircum{}\{2 pi i/3\} r, so r carries a SHEET INDEX in Z\_3 that advances by one per passage through the r=0 branch point and returns after three -- exactly a crossing count's transformation law, and derived in the corpus from cosmic time's imaginary period before any of this. Closure then MEANS returning to the same branch rather than stipulating a rule [borrowed from P3 / P14]). 
\emph{Computes:} the criterion tabulated, the incidence exhausted over all 9 pairs, the leg factorisation enumerated, the sheet arithmetic exact
 \emph{Bound:} prop:wall's mode has NOT been carried through r=0 to read off its sheet advance (L-78) -- one computation on an existing solution. READING A (hinge = colour) is forced by the causal trichotomy, so P14's 'three walls = three generations' must be revised and generation sought on the turnaround three (L-67, reaching the dimension result). And 15 = 12 + 3 is REGISTERED NOT ASSERTED, flagged as the most numerology-exposed claim in the reading, refutable by the three wall states carrying a horn label (L-80, unchecked)
\ \emph{Run:} \texttt{python3 receipts/P03\_SdS\_slicing/P03\_sheet\_index\_and\_factors.py}
\item[\rcptlabel{P03_wall_is_the_graze_point}{P3R46}\textbf{P3R46}\quad\texttt{P03\_wall\_is\_the\_graze\_point}]\hfill\textsf{[OK]}\\
\textit{sec:tour / sec:chirality} \ (THE MATTER SECTOR'S WALL IDENTIFIED WITH THE THROAT POINT WHERE THE NULL LEGS GRAZE -- two independent derivations, one set of three points, never before connected. P14 reaches the walls from the signed radius ('r=0 is NOT the throat's CENTRE but a point ON the throat circle, the back, fixed relative to the hinge; the Z\_3-fixed centre carries NO wall; the walls lie at polar 180/300/60, a Z\_3-orbit, EACH ANTIPODAL TO ITS HINGE'). Receipt P03\_twelve\_null\_legs reaches 60/180/300 from null tangency. VERIFIED IDENTICAL, and the pairing is exact: each wall is hinge k's back AND is where the OTHER TWO hinges' null tangents touch the hole -- so the hinge<->wall pairing is the ANTIPODAL map, canonical rather than handedness-dependent. AND a hinge's own two null legs cross the other two hinges' walls and NEVER its own (3 of 3), so the 1+2 appears a third time, now on walls. CONSEQUENCE: since the graze points ARE the walls and a wall is where prop:wall binds a chiral zero-mode by Jackiw-Rebbi, a winding is a WALL-CROSSING COUNT -- the natural label for a wall-bound field rather than an import; one lap = three wall crossings = the whole Z\_3 orbit, each wall being an r=0 branch point [borrowed from P3 / P14]). 
\emph{Computes:} both point sets constructed independently and compared; the pairing and the never-its-own property exhausted over all three hinges
 \emph{Bound:} what this file establishes is that L-74's antecedent is the corpus's own structure; the step it left owed -- that the zero-mode's label transforms as a crossing count -- was RUN at c54.25 (L-78), the $\sqrt{f}$ cancelling exactly in $W\,d\ell=\lambda\,dr/r$ so that $\psi\sim r^{-\lambda}$ and $\psi\mapsto\omega^{-\lambda}\psi$ per crossing, i.e. triality $=-\lambda\bmod3$; L-72 and L-74 are closed (r2886). AND THE WELD CONFLICT THIS RECORDED IS RESOLVED RATHER THAN STANDING: P14's sec:whichthree reseats the generations off the walls onto the turnaround's deck $\mathbb{Z}_3$, so the hinge three is the WITHIN-STATE index and the two threes are related by no covering construction -- walls and hinges are not asked to be both, and L-67 is discharged by that reseating
\ \emph{Run:} \texttt{python3 receipts/P03\_SdS\_slicing/P03\_wall\_is\_the\_graze\_point.py}
\item[\rcptlabel{P03_winding_vs_triality_closure}{P3R59}\textbf{P3R59}\quad\texttt{P03\_winding\_vs\_triality\_closure}]\hfill\textsf{[OK]}\\
\textit{sec:winding / P14 sec:chirality} \ (THE TWO ADMISSION RULES ARE INDEPENDENT CONDITIONS, AND THE ELEVEN CHANNELS CANNOT SEPARATE THEM. The corpus carries (W) P3's total-winding rule, sum n\_i = 0 mod 3, a statement about the COVERING of r; and (T) P14's total-residue rule, sum t\_i = 0 mod 3, with t = -lambda mod 3 read off the ANGULAR problem (L-78). Both return 11/11 on the channel set. BUT every constituent in that set is assigned n = t at the outset -- quark (+1,+1), antiquark (-1,-1) -- so sum n = sum t identically on all eleven and the set returns one verdict for both rules whatever the physics. ** So the 11/11 tests the admission rule and NOT the identification. ** Separating them requires a constituent with n != t mod 3; the independence is then exhaustive at sizes 2-3, 180 configurations admitted by (W) and refused by (T) and 180 the other way. CONSEQUENCE: 'the identification of the triality class with the charge's fractional part remains the assumption it was' is exactly the per-constituent condition n = t mod 3, and without it the two rules are logically independent). 
\emph{Computes:} both rules on all eleven channels; the sum n = sum t identity over the set; the smallest separating configuration; the two-way independence count exhausted over sizes 2--3
 \emph{Bound:} NOT CLAIMED that the geometry does or does not admit a constituent with n != t -- that is the open condition, not a result here. NOT CLAIMED that the 11/11 is weakened: it stands for the rule it tests. Takes up exactly what \texttt{B24\_the\_triality\_test\_run} marks as untouched (``the derived-vs-assumed identification''), and does not duplicate it -- B24 tests t against the configuration count, this tests (W) against (T)
\ \emph{Run:} \texttt{python3 receipts/P03\_SdS\_slicing/P03\_winding\_vs\_triality\_closure.py}
\item[\rcptlabel{P03_closure_under_the_three_vantage_reading}{P3R60}\textbf{P3R60}\quad\texttt{P03\_closure\_under\_the\_three\_vantage\_reading}]\hfill\textsf{[OK]}\\
\textit{sec:winding / P14 sec:chirality} \ (THE COMPOSITE ADMISSION RULE IS NOT FORK-INSENSITIVE, AND IT IS THE ONE ENTRY P14R28's DECISION TABLE DOES NOT CARRY. That table checked triality, the thirds, the three classes and route-dependence against the (i)/(ii) fork on the signed radius and found all four insensitive. Under (ii) -- three radii, a crossing branching only the owner's vantage -- a circuit crossing wall W with multiplicity m\_W has monodromy diag(omega\^(-lambda m\_A), omega\^(-lambda m\_B), omega\^(-lambda m\_C)), so IDENTITY requires lambda m\_W = 0 mod 3 for EACH wall separately, where P3's rule conditions on the TOTAL. Exhaustive over multiplicities 0--3 and lambda 1,2: 28 configurations separate them, the smallest being THE LAP ITSELF -- m=(1,1,1), total 3, admitted by the total rule for every lambda and trivial only at lambda = 0 mod 3. ** That is P14R28's own headline read from the other side: 'under (ii) the lap acts by omega\^(-lambda), and only triality zero returns.' ** Implication computed both ways: (ii)-triviality ==> sum n = 0 mod 3 TRUE, converse FALSE -- so the total-crossing rule is NECESSARY and not sufficient, and is the shadow on the total of a condition that is per-wall). 
\emph{Computes:} the monodromy under (ii) for all multiplicities 0--3 at lambda = 1,2; the separating set; the lap at three values of lambda; necessity and sufficiency each exhausted
 \emph{Bound:} NOT CLAIMED that P3's rule is wrong -- it is correct under (i) and a correct necessary condition under (ii), and every configuration the corpus has tested is admitted or refused identically by both, the eleven being stated in constituent content and not resolving m\_W. NOT CLAIMED that the fork should be reopened; (ii) stands on the sentence P14R28 used to decide it. The single-mode circuit is the object computed; a composite of constituents at different lambda is the per-constituent version of the same condition and is not separately enumerated here
\ \emph{Run:} \texttt{python3 receipts/P03\_SdS\_slicing/P03\_closure\_under\_the\_three\_vantage\_reading.py}
\item[\rcptlabel{P03_the_operative_crossing_count}{P3R61}\textbf{P3R61}\quad\texttt{P03\_the\_operative\_crossing\_count}]\hfill\textsf{[OK]}\\
\textit{sec:winding / P14 sec:chirality} \ (THE COUNT THAT ACTS ON A BOUND MODE IS NOT THE COUNT THE WINDING MEASURES, AND THE TWO ARE LOCKED. Under reading (ii) a crossing of wall W branches only W's OWNER, so for a mode bound in vantage k the crossings that act on its own component are the crossings of WALL k alone. Enumerated over all twelve routes between punctures at distinct hinges, with hinges at 0/120/240 and wall\_j antipodal at (120j+180): the realised pairs are (n,nu) = (1,0) and (2,1) only, so ** nu = n - 1 on every route **, forced by the antipodal map -- wall\_m lies between hinges i and j whenever \{i,j,m\}=\{0,1,2\}, so the SHORT route from hinge k crosses only the third hinge's wall and the LONG route crosses hinge k's own wall and the destination's. CONSEQUENCE FOR THE OPEN CONDITION: the own-slot monodromy is omega\^(t\*nu) = omega\^(t(n-1)), so n and t enter as a PRODUCT and the identification n = t mod 3 is NOT what the geometry imposes. A route with n=1 leaves the mode's own component untouched whatever its triality; a route with n=2 multiplies it by omega\^t). 
\emph{Computes:} the twelve inter-hinge routes with both counts and the own-slot monodromy; the realised (n,nu) set; the relation checked route by route; the own-slot monodromy tabulated over n and t
 \emph{Bound:} nu = n - 1 is computed on routes between punctures at DISTINCT hinges, which is the domain P3 sec:winding defines a winding on, and is NOT a general identity -- a full lap has n=3 and nu=1, and the trivial route n=0 and nu=0. NOT DERIVED: the identification n = t mod 3, which remains open. What the binding fixes is WHICH crossings act (own-wall only); what the geometry fixes is nu = n-1; the triality is set by the angular problem and nothing here relates it to either
\ \emph{Run:} \texttt{python3 receipts/P03\_SdS\_slicing/P03\_the\_operative\_crossing\_count.py}
\item[\rcptlabel{P03_can_the_charge_be_read_from_the_acting_count}{P3R62}\textbf{P3R62}\quad\texttt{P03\_can\_the\_charge\_be\_read\_from\_the\_acting\_count}]\hfill\textsf{[OK]}\\
\textit{sec:winding} \ (NO -- AND THE REASON IS THE THIRDS THEMSELVES. P3R49 gave two counts on a route: n, the signed graze-point count the winding measures, and nu, the own-wall count that actually acts on the bound mode. Since nu is operative, the natural next reading is that the charge should be read from it; run through P3's own derivation -- partner relation, three constituents, closure of both horn-headed triples -- that reading FAILS. charge = n/3: the two routes carry +1/3 and -2/3, the partner relation y = x-1 holds, both closure totals are integers and the thirds are FORCED. charge = nu/3: the routes carry 0 and -1/3, the partner relation FAILS (difference -1/3, not one lap). charge = nu: the partner relation holds but nu is integer-valued by construction, closure is satisfied identically and NO denominator is forced. ** So the thirds exist precisely BECAUSE the charge-carrying count and the monodromy-acting count are different counts -- one running over all three walls and fractional in thirds of a lap, the other over one wall and integral. ** Signed form of P3R49's relation also checked: nu = n - sgn(n)). 
\emph{Computes:} the signed two-count table; P3's derivation run on all three candidate labels with the partner relation and both closure totals evaluated; the signed relation
 \emph{Bound:} this CLOSES one candidate route to deriving the standing identification and does not bear on the identification itself -- n = t mod 3 relates the charge-carrying count to the triality and is untouched here. NOT CLAIMED that P3's derivation needed defending: it is unchanged, and what is added is why the obvious substitution into it does not work
\ \emph{Run:} \texttt{python3 receipts/P03\_SdS\_slicing/P03\_can\_the\_charge\_be\_read\_from\_the\_acting\_count.py}
\end{description}\endgroup
